%!TEX root = ../thesis.tex

\chapter[Conclusion and Future Directions]{Conclusion and Future Directions}
\label{chapter:conclusion}

This thesis has endeavored to push the frontiers of sampling by introducing unifying theoretical frameworks and scalable algorithmic solutions across generative modelling, conditional sampling, and large-scale Bayesian inference. We first developed a novel variational perspective that successfully bridged the traditionally distinct domains of variational inference, optimal transport, and diffusion modelling, leveraging novel theorems within stochastic calculus. This allowed us to demonstrate that many contemporary diffusion-based algorithms for both generative modelling and sampling emerge as specific instances of our proposed framework. Subsequently, we established a unifying mathematical lens for conditional diffusion methods using Doob's \(h\)-transform. Finally, we pivoted to the pressing challenge of large-scale Bayesian inference, pioneering efficient stochastic gradient descent (SGD)-based methodologies. While these contributions have addressed key challenges in unifying diverse sampling paradigms and in scaling sophisticated Bayesian methods to contemporary problem sizes, the journey also illuminates numerous exciting avenues for future exploration.

The development of our unifying variational framework in \Cref{sec:path_unifying_framework} for generative modelling and sampling culminated in several key theoretical and practical advances. Within this framework, we were able to unify a large number of algorithms in \Cref{sec:connections_to_existing_methodologies}, including score matching, ergodic samplers, finite-hitting samplers, optimal transport algorithms, action matching, and diffusion normalising flows, to name some. We highlighted links to foundational fluctuation theorems from statistical physics like those of Crooks and Jarzynski in \Cref{prop:controlled_crooks_fluctuation_theorem}. The practical culmination of this framework was the Controlled Monte Carlo Diffusions (CMCD) inference scheme in \Cref{sec:controlled_monte_carlo_diffusions}, which achieved state-of-the-art results across a suite of challenging inference benchmarks (see \Cref{sec:cmcd_experiments}). We believe this experimental success is partly attributable to CMCD striking an effective balance between parametrising a flexible family of distributions and being sufficiently constrained to make learning the sampler tractable and not overly expensive, a crucial consideration highlighted by prior theoretical work. Looking ahead, promising avenues stemming from this work include the exploration of optimal annealing schedules for the path integral \(\pi_t\), potentially adapting strategies from adaptive importance sampling \parencite{goshtasbpour2023adaptive}, and the investigation of alternative path-space divergences which might offer different trade-offs in convergence or expressivity \parencite{nusken2021solving, richter2023improved, midgley2022flow}.

Shifting focus to the critical task of conditional sampling, we introduced a unifying mathematical perspective based on Doob's $h$-transform in \Cref{chapter:deft}. This allowed for a clearer understanding and classification of diverse conditional diffusion methods in \Cref{sec:generalised_htrans}. From this theoretical standpoint, we proposed Doob’s \(h\)-transform Efficient Fine-Tuning (DEFT), a novel parameter-efficient conditional fine-tuning method in \Cref{sec:training_objectives}. A key advantage of DEFT is its ability to operate without requiring backpropagation through large, pre-trained score networks during the fine-tuning or sampling phases, resulting in significant efficiency gains. Our evaluations on several image reconstruction tasks demonstrated that DEFT reliably outperformed standard methods in terms of sampling time, reconstruction quality, and perceptual similarity metrics (see \Cref{sec:image_rec}). Notably, while DEFT necessitates an additional training phase on a small dataset of paired measurements, its overall speed often surpasses existing baselines due to its capacity for fewer sampling steps during evaluation and the absence of backpropagation at inference.
Despite the strengths of DEFT, certain limitations point towards important future research directions. The current DEFT framework relies on the availability of a (small) fine-tuning dataset, which contrasts with zero-shot conditional sampling approaches such as DPS or \(\Pi\)GDM that operate without paired data. Fine-tuning, even on small datasets, carries an inherent risk of overfitting or reproducing biases present in that specific data. Furthermore, unlike some zero-shot methods, DEFT in its current form does not explicitly assume knowledge of the forward operator. Future iterations could investigate incorporating the forward operator as an inductive bias within the network architecture to potentially enhance performance and data efficiency. We also proposed a preliminary zero-shot approach via a stochastic control loss, which requires only a single observation to learn the \(h\)-transform. While initial results on MNIST were promising (see \Cref{sec:stoch_cont_exp}), the computational burden of simulating the full SDE at each iteration currently limits its feasibility for high-dimensional data. However, recent promising work on partial trajectory optimisation for diffusion models \citep{domingo2024adjoint} may significantly reduce this computational cost, making such stochastic control objectives competitive with existing methods and a rich area for further development. Finally, we extend our method in \textcite{denker2025iterative}, a recent follow-up that iteratively refines the $h$-transform using a synthetic dataset resampled with path-based importance weights, showing promising results for dataset-free finetuning using the $h$-transform and the pathwise RND.

The third major section of this thesis tackled the scalability of sampling from Bayesian posteriors, particularly through the lens of stochastic optimisation. For Gaussian Processes (GPs), we introduced stochastic dual descent in \Cref{sec:sgd_for_sampling}, a specialised first-order stochastic optimisation algorithm tailored for computing GP predictions and, crucially, for drawing posterior samples. By combining insights from the broader optimisation literature with GP-specific ablations, we arrived at an algorithm that is simultaneously simple to implement and robustly matches or exceeds the performance of relevant contemporary baselines. We demonstrated that stochastic dual descent yields strong performance on standard regression benchmarks, excels in a large-scale Bayesian optimisation benchmark where sample quality is paramount, and even matches the performance of state-of-the-art graph neural networks on a challenging molecular binding affinity prediction task. Finally, we used our stochastic sampling algorithms to propose a novel \emph{warm-start pathwise} algorithm for stochastic trace estimation in GP hyperparameter optimisation in \Cref{sec:pathwise_gradients_hparam}, resulting in speed-ups of over $72\times$ compared to the standard estimator in standard linear solvers.

Building on these insights into scalable sampling for structured Gaussian models, we then addressed the complexities of Bayesian Neural Networks (BNNs) in \Cref{chapter:bnn_sampling}. Our work introduced a sample-based approximation to both inference and hyperparameter selection in Gaussian linear multi-output models, with a particular focus on scaling the Linearised Laplace Approximation (LLA). This approach proved to be asymptotically unbiased and, critically, allowed us to scale the LLA method to large ResNet models on CIFAR100 and even ImageNet without discarding crucial covariance information—a step that was computationally intractable with previous methods (see \Cref{sec:bnn_experiments}). The efficacy of this sample-based EM approach was further demonstrated on a high-resolution tomographic reconstruction task in \Cref{sec:DIP}, where it decreased the cost of hyperparameter selection by two orders of magnitude. Importantly, the uncertainty estimates obtained through our method were shown to be well-calibrated not just marginally, but also jointly across predictions.

The advancements in scalable Bayesian inference presented open up numerous possibilities. The ability to obtain well-calibrated joint posterior samples efficiently for large models is of particular interest in fields such as active learning and reinforcement learning, where decisions often depend on understanding correlated uncertainties and where posterior sampling is frequently required. Future work in this domain could focus on extending these SGD-based sampling and hyperparameter optimisation techniques to even more complex, non-conjugate Bayesian models, potentially by integrating them with more sophisticated variational approximations or advanced MCMC techniques designed for non-Gaussian posteriors. Further research into adaptive step-size schemes and preconditioners specifically designed for these sampling-via-optimisation objectives could also yield substantial performance gains.

In conclusion, this thesis has contributed novel unifying frameworks that bring clarity to the rapidly evolving landscape of generative modelling and conditional sampling, while also delivering practical, scalable algorithms for some of the most challenging problems in Bayesian inference. From establishing connections between variational inference, optimal transport, and diffusion models, to providing a principled understanding of conditional sampling via Doob's transform, and finally to leveraging stochastic optimisation for massive-scale Bayesian computation, we have aimed to provide both theoretical insight and algorithmic innovation. The path forward remains rich with opportunities to refine these methods, explore new connections between disparate fields, and ultimately, to continue enhancing our ability to learn from data and quantify uncertainty in an increasingly complex world.
% \paragraph*{Limitations and future work} The DEFT framework uses
% % relies on the availability of 
% a (small) fine-tuning dataset, 
% % which is 
% in contrast to zero-shot conditional sampling approaches, e.g., DPS \cite{chung2022diffusion} or $\Pi$GDM \cite{song2022pseudoinverse}. Fine-tuning on small datasets may have the risk of 
% % reproducing and 
% overfitting to biases inherent in the data. 
% In contrast to zero-shot conditional sampling, DEFT assumes no knowledge of the forward operator. However, the forward operator can be incorporated as an inductive bias within the network architecture to improve performance.
% % 
% We also proposed a zero-shot approach through the optimal control loss in \Cref{sec:stochastic_control}, which only needs a single observation $\m$ to learn the $h$-transform. Though we show promising results scaling this approach to the MNIST dataset, the computational burden of simulating the full SDE at each iteration is still high, which might make this optimal control loss infeasible for high-dimensional data. However, there is promising recent work on partial trajectory optimisation \cite{zhang2023diffusion}, which may reduce the computational burden of the stochastic control objective, making it competitive with existing methods.